\section{Tabular Graph \texorpdfstring{\emojitabtitle}{Tabular Graph}}
\label{sec-tab}

Tabular data is another type of structured data that is widely used in real-world applications~\citep{sahakyan2021explainable} and is typically stored in relational databases~\citep{codd2007relational}. Tabular data may consist of a single table containing samples and their attributes, or multiple tables that share primary and foreign keys. LLMs have been explored to process and solve tasks involving tabular data, primarily by transforming the tabular data into text through serialization~\citep{fang2024large, singha2023tabular}. However, such serialization may lead to some issues: (1) In the case of a single table, the feature columns should exhibit permutation invariance, but serializing the table data may disrupt this invariance; (2) For multiple tables, one table may be connected to another through primary/foreign keys, and leveraging these relationships is crucial for various tasks. Therefore, graph structures can be a suitable representation for tabular data. Additionally, when tables contain too many rows to fit within the context window of LLMs, graphs can facilitate efficient retrieval.

\subsection{Application Tasks}
Tables are widely used in real-world scenarios to store features and relationships between data. Understanding the structure of tabular data is crucial. As a result, there are numerous tasks that can benefit from the use of tabular graphs and below we list a few representative ones. 

\begin{itemize}[leftmargin=*]
    \item \textbf{Node-level tasks}: Node-level tasks include node classification and node regression, applied to tasks like cell type prediction~\citep{jin2023tabprompt}, fraud detection~\citep{rao2020xfraud, singh2023graphfc}, outlier detection~\citep{goodge2022lunar}, and click-through rate (CTR) prediction~\citep{li2019fi, du2022learning}.

    \item \textbf{Link-level tasks}: Link-level tasks involve link prediction, edge classification, and edge regression. Many tabular data tasks can be modeled as link-level tasks, such as data imputation~\citep{zhong2023data, you2020handling} and recommendation~\citep{robinson2024relbench}.

    \item \textbf{Graph-level tasks}: Graph-level tasks aim to predict the properties of the entire graph, such as table type classification and table similarity prediction~\citep{wang2021tuta, jia2023getpt, jin2023tabprompt}.

    \item \textbf{Table question answering}: Table QA involves generating answers by understanding and reasoning over tabular data. This task requires comprehension of both the content and their relationships within tables, making graph structures suitable for encoding such information. For example, \citet{zhang2020cfgnn, zhang2020graph} utilize graphs to enhance Table QA.

    \item \textbf{Table retrieval}: Table retrieval focuses on retrieving semantically relevant tables based on natural language queries~\citep{wang2021retrieving, cheng2021hitab}.
\end{itemize}

\subsection{Tabular Graph Construction}
Graphs are used in tabular data learning to model high-order feature interactions, high-order instance relationships, and relationships between instances across multiple tables. There are typically two types of nodes: instance nodes, which represent each row of a table, and feature nodes, which represent individual features. Generally, the following graphs are constructed:

\begin{itemize}[leftmargin=*]
    \item \textbf{Instance Graph}: An instance graph connects a table's rows (instances), modeling the relationships between instances. It is particularly useful for retrieving relevant instances within tabular data. There are mainly two approaches for constructing instance graphs:
    
    \begin{itemize}[leftmargin=*]
        \item \textit{Rule-based methods}: Instances are connected based on predefined rules. For example, heuristics derived from expert knowledge can be used to connect instances that share certain features~\citep{lu2021weighted}. Expert knowledge can be leveraged for the graph construction~\citep{parisot2018disease}. Another common approach is similarity-based methods, such as connecting instances through K-Nearest Neighbors (KNN)~\citep{gu2022structure, errica2024class} or based on similarity exceeding a threshold~\citep{spinelli2020missing, chen2023gedi}.
        \item \textit{Learnable methods}: In this approach, an instance graph is typically initialized using rules or heuristics. The edge weights are then adjusted during the learning process, allowing the model to dynamically refine the graph structure over time~\citep{kang2021k, liao2023tabgsl, cosmo2020latent, kazi2022differentiable}.
    \end{itemize}
    

    \item \textbf{Feature Graph:} A feature graph connects features, with edges representing correlations between pairs of features. Typically, feature relationships are modeled in a learnable manner~\citep{yan2023t2g, zhou2022table2graph}. Some works construct the feature graph by leveraging feature similarity~\citep{guo2021tabgnn, li2019fi}, while others use heuristics based on expert knowledge to establish connections~\citep{li2024graph}. Additionally, certain approaches link features if they belong to the same instance~\citep{rao2020xfraud}.

    \item \textbf{Instance-Feature graphs:} The instance-feature graph is a heterogeneous graph, which connects the instance with their corresponding features~\citep{guo2021dual, you2020handling, wang2021retrieving, wu2021towards}. 

    \item \textbf{Cell Graph:} Cell graphs treat each cell in the table as a node. \citet{xue2021tgrnet} build a cell graph, where each cell node contains both spatial and logical locations attributes while the adjacency matrix represents the neighbor relation or the same-row and same-column relation between two cells.

    \item \textbf{Tabular Hypergraph:} A hypergraph is a generalization of a graph in which an edge, known as a hyperedge, can connect multiple nodes. Since tables are often invariant to arbitrary row and column permutations, a hypergraph is a suitable structure for modeling tabular data~\citep{zahradnik2023deep}. Specifically, \citet{chen2024hytrel} model each row and column in a table as a hyperedge for pretraining purposes. Additionally, \citet{du2022learning} construct a hypergraph from tabular data to capture relationships between instances effectively.

    Additionally, \citet{cvetkov2023relational} model table as a knowledge graph, where the feature values or rows represent nodes, and the column names represent the relations. \citet{zhang2020graph} build a heterogeneous graph based on various predefined relations. \citet{wang2021tuta} construct trees for non-relational tables based on the hierarchical relations. 

    Previous approaches primarily focus on constructing graphs based on a single table. However, in relational databases, there are often multiple tables. A common method is to first merge these tables into a single table, and then apply the previously mentioned methods to the merged table~\citep{kanter2015deep}. However, this approach relies on manually joining the tables as part of a feature engineering step, which requires significant effort and substantial domain expertise. In the following, we will introduce cross-table graphs, which directly build a graph from multiple tables without the need for manual merging.

    \item \textbf{Cross-table Graphs:} Cross-table graphs for tabular data typically connect different instances across multiple tables, where nodes usually represent rows in each table and edges are defined by primary-foreign key relationships~\citep{fey2023relational, cvitkovic2020supervised, zhang2023gfs}. Additionally,\citet{wang20244dbinfer} propose Row2N/E: if a table has a primary key (PK), each row is treated as a node; however, if there are also two foreign key (FK) columns, each row is instead represented as an edge. These graphs model relationships across multiple tables, enabling the integration of information from different perspectives to facilitate various tasks.

    Additionally, \citet{bai2021atj} build a hypergraph based on multiple tables, where the primary or foreign key as nodes, and the nodes within the same table as hyperedge.

\end{itemize}


\subsection{Retriever}
Modeling tabular data into graphs is still in its early stage, and most works follow popular retrieval methods, as introduced in Sections~\ref{sec:retriver}. Additionally, \citet{fey2023relational} retrieve subgraphs based on timestamps to construct time-consistent computational graphs, while \citet{cvitkovic2020supervised} adopt a deterministic heuristic to select subgraphs.

\subsection{Generator}
Most existing methods that leverage tabular graphs still rely on GNNs or Graph Transformers as generators~\citep{li2024graph}. Besides, \citet{wang20244dbinfer} fuse both GNNs and tabular based predictors, such as DeepFM~\citep{guo2017deepfm}, FT-Transformer~\citep{gorishniy2021revisiting}, XGBoost~\citep{chen2016xgboost} and AutoGluon~\citep{erickson2020autogluon}. \citet{ivanov2021boost} jointly train gradient boosted decision tree (GBDT) and GNN. Recent advancements have seen the application of LLMs to tabular data tasks. This is typically achieved by converting tabular data into sequential formats suitable for LLM input, as explored by \citet{fang2024large}, \citet{dong2024large}, and \citet{sui2024table}. However, this transformation can lead to the loss of inherent structural information present in the original tabular format. With advancements in the application of LLMs to graph data, LLMs can be further enhanced to handle various tasks more effectively with the support of tabular graphs.

\subsection{Resources and Tools}
In this section, we list some data sources and tools for GraphRAG on tabular graphs.

\subsubsection{Data Resources}
Relational machine learning has recently gained popularity, leading to the development of several benchmarks and datasets for tabular data. For example:

\begin{itemize}[leftmargin=*]
    \item \textbf{RelBench}~\citep{robinson2024relbench, fey2023relational}\footnote{https://relbench.stanford.edu/} is a collection of realistic, large-scale, and diverse benchmark datasets for machine learning on relational databases. It includes various tasks, such as Node-level prediction tasks (e.g., multi-class classification, multi-label classification, regression), Link prediction tasks, Temporal and static prediction tasks. The benchmark features several datasets, including rel-amazon, rel-stack, rel-trial, rel-f1, rel-hm, rel-event, and rel-avito, providing a robust platform for evaluating models across multiple relational data scenarios and tasks.
    
    \item \textbf{TabGraphs}~\citep{bazhenov2024tabgraphs} evaluates various models — including simple baselines, tabular models, and GNNs — on graphs with tabular node features. This benchmark includes several TabGraphs datasets: tolokers-tab, questions-tab, city-reviews, browser-games, hm-categories, hm-prices, city-roads-M, city-roads-L, avazu-devices, web-fraud and web-traffic. 
    \item \textbf{Tabular-benchmark}~\citep{grinsztajn2022tree} \footnote{https://github.com/LeoGrin/tabular-benchmark} benchmarks 45 tabular datasets from diverse domains, comparing the performance of several tree-based and deep learning models. They evaluate tasks such as numerical classification, numerical regression, categorical classification, and categorical regression. 
    \item \citet{shwartz2022tabular} benchmark 11 tabular datasets, such as 
    MSLR~\citep{qin2013introducing},  Forest Cover Type, Higgs Boson, and Year Prediction~\citep{dua2017uci}.
    evaluating several tree-based models, deep learning models, and ensemble models.
    \item \textbf{DBInfer Benchmark}~\citep{wang20244dbinfer}\footnote{https://github.com/awslabs/multi-table-benchmark}  is a set of benchmarks for measuring machine learning solutions over data stored as multiple tables. It includes several large-scale relational database (RDB) datasets—such as AVS, OB, DN, RR, AB, SE, MAG, and SE—with tasks like retention, CTR, purchase, CVR, churn, rating, popularity, venue prediction, citation, charge, and prepay. The benchmark provides implementations of various baselines, including popular tabular models with and without auto-feature-engineering methods, as well as Graph Neural Networks, making it a robust resource for multi-table learning evaluations.
    \item \textbf{RTDL}\footnote{https://github.com/yandex-research/rtdl}: RTDL (Research on Tabular Deep Learning) is a collection of papers and packages on deep learning for tabular data. It provides several popular tabular deep learning models.
    \item \textbf{OpenML}~\footnote{https://www.openml.org/} is an open platform for sharing datasets, algorithms, and experiments. It provides a large collection of machine learning datasets across various domains, including many tabular datasets.
    \item \textbf{Kaggle}~\footnote{https://www.kaggle.com/} hosts a wide array of datasets for data science competitions, making it a valuable resource for benchmarking tabular machine learning models on diverse real-world data.
    \item \textbf{UC Irvine Machine Learning Repository}~\citep{dua2017uci}\footnote{https://archive.ics.uci.edu/} is a collection of databases, domain theories, and data generators that are used by the machine learning community for the empirical analysis of machine learning algorithms.
    \item \textbf{HiTab}~\citep{cheng2021hitab} is a hierarchical table dataset for question answering and natural language generation.
\end{itemize}


\subsubsection{Tools}
In this subsection, we list several tools that provide essential functionalities for data preparation, model training, and evaluation.  These tools can be effectively leveraged for GraphRAG applications on tabular graphs.

\begin{itemize}[leftmargin=*]
    \item \textbf{PyTorch Tabular}~\citep{joseph2021pytorch}\footnote{https://github.com/manujosephv/pytorch\_tabular}: PyTorch Tabular is a powerful library built on PyTorch designed to simplify the application of deep learning techniques to tabular data. It provides several data preprocessing functions, including normalization, standardization, encoding of categorical features, and dataloader preparation. Additionally, it includes a variety of tabular machine-learning models and evaluation functions, making it a comprehensive tool for handling tabular data.
    \item \textbf{DeepTables}~\citep{deeptables}\footnote{https://github.com/DataCanvasIO/DeepTables}: DeepTables is an easy-to-use toolkit that harnesses the power of deep learning for tabular data. Built on TensorFlow, it is designed to address classification and regression tasks on tabular data. DeepTables offers a variety of tabular models and supports AutoML.
    \item \textbf{PyTorch Frame}~\citep{hu2024pytorch}\footnote{https://github.com/pyg-team/pytorch-frame}: PyTorch Frame is a deep learning extension for PyTorch, designed for heterogeneous tabular data with different column types, including numerical, categorical, time, text, and images. t provides a wide range of tabular models and supports integration with diverse model architectures, including Large Language Models.
\end{itemize}

